\section{Simulation and resampling}

Three ways to generate outcomes are distinguished here: replaying the past (historical simulation), resampling the observed data (bootstrap) and drawing from a specified probability model (Monte Carlo). The risk-measure applications, filtered historical simulation and dynamic models are owned by \hyperref[meth:var-es-estimation-methods]{VaR and ES estimation}.

\begin{qfmethod}[Historical simulation: replaying past changes on today's position, equal and age weights, and what the method assumes]\label{meth:historical-simulation}
\qffield{Statement} Historical simulation ``involves using past data as a guide to what will happen in the future''. With data on the most recent 501 days, each day's percentage changes in the market variables define one of 500 scenarios; scenario $i$ sets tomorrow's value of a variable to $v_n v_i/v_{i-1}$, where $v_n$ is today's value. Revaluing the current portfolio under each scenario gives a loss distribution. Hull estimates its 99th percentile, the one-day 99 percent VaR, as the fifth-highest of 500 losses, noting that ``a case can be made for using the fifth-highest loss, the sixth-highest loss, or an average of the two''; the estimate holds if the past 501 days are ``representative of what will happen'' \citep[chap.~22, sec.~22.2, p.~517 and n.~2]{s7-hulloptions}. Expected shortfall averages the losses in the tail, the five worst in his example (p.~520). To give recent scenarios more weight, scenario $i$ can receive $\lambda^{n-i}(1-\lambda)/(1-\lambda^n)$, with $0<\lambda<1$; the 99 percent VaR is then ``the loss for the first scenario which is such that the cumulative weight exceeds 0.01'', and ES weights the losses above it by their weights and the VaR loss by the weight remaining to make up 0.01 (p.~521).

\qffield{Notation} $n$ scenarios; losses $L_1,\dots,L_n$ of the \emph{current} portfolio under each historical change, with weights $w_i$ summing to one ($1/n$ for equal weights); $p=1-c$. This compendium uses the lower loss quantile and exact-tail-mass ES of \hyperref[meth:var-es-estimation-methods]{VaR and ES estimation} \citep[Definitions~2.2 and~2.6, pp.~3 and~5; Proposition~3.2, p.~6]{acerbi-tasche-2001}. Rank the losses from the largest, $\ell_{(1)}\ge\ell_{(2)}\ge\dots$, merging ties, and let $W_k$ be the weight of the $k$ largest. Interpretation: $\mathrm{VaR}_c=\ell_{(k)}$ for the first $k$ with $W_k>p$, which is Hull's weighted rule, and
\[ \mathrm{ES}_c=\frac1p\Big(\sum_{j<k}w_{(j)}\ell_{(j)}+(p-W_{k-1})\,\ell_{(k)}\Big), \]
the average over a tail of mass exactly $p$, taking only part of the weight at the VaR loss when needed. With equal weights and $np$ an integer, VaR is the $(np+1)$th largest loss, one rank below Hull's equal-weight choice of the $np$th, and ES is the mean of the $np$ largest; at $np=2.5$, $\mathrm{ES}=(\ell_{(1)}+\ell_{(2)}+0.5\,\ell_{(3)})/2.5$. The two VaR conventions differ only in which side of a jump in the empirical distribution is reported; ES is the same under both.

\qffield{Conditions} The method keeps the joint behaviour and fat tails of the sample but can produce no loss larger than its worst day scaled to today's position. It assumes the window is representative: a calm window understates risk, a window containing a crisis states it until the crisis drops out. Age weighting reacts faster but uses fewer effective observations. Using returns or price ratios of today's holdings, not past portfolio P\&L, keeps the position current (Interpretation).

\qffield{Practice} Example, original: a portfolio of 1{,}000{,}000, 60 percent in asset~A and 40 percent in asset~B, and 20 invented daily return pairs. Revalued under each day, the worst losses are 18{,}600 (day~7) and 12{,}800 (day~14). The next are 8{,}400 (day~2) and 7{,}600 (day~19). At 90 percent, $p=0.1$ and $np=2$. With equal weights the two worst losses carry exactly the tail mass, $W_2=0.10$, which does not exceed $p$, so VaR is the third-ranked loss, 8{,}400, and ES the average of the two worst, 15{,}700. Hull's equal-weight rank would report the second-worst loss, 12{,}800, as VaR; the ES is unchanged. With $\lambda=0.9$ the two worst days, old in the window, carry weights 0.029 and 0.060, $W_2=0.089$, and the third-ranked loss (weight 0.017) takes the cumulative weight past 0.10, so VaR is again 8{,}400. ES becomes $(0.029\times18{,}600+0.060\times12{,}800+0.011\times8{,}400)/0.1\approx14{,}000$: down-weighting the old losses lowers ES by about 11 percent, while VaR does not move on this sample. Comparing a weighted VaR with an equal-weight VaR computed under another rank convention would wrongly attribute the convention's effect to the weighting.

\qffield{Pitfalls} Reporting a historical VaR without its window, weighting and quantile convention, or comparing estimates computed under different conventions. Computing ES as the mean of losses strictly beyond VaR. At the VaR of 8{,}400 that average happens to equal the ES here, 15{,}700, because the two losses above it carry exactly the tail mass; it omits the boundary weight whenever the losses above VaR carry less than $p$, as under the age weights. Averaging only the losses beyond Hull's equal-weight VaR of 12{,}800 would instead leave the 18{,}600 loss alone. Expecting a loss beyond the worst day in the window. Mixing scenario returns with the past composition of the portfolio.

\qffield{See also} \hyperref[meth:var-es-estimation-methods]{VaR and ES estimation} (filtered HS, backtests); \hyperref[def:return-sample-statistics]{percentile conventions}.

\end{qfmethod}

\begin{qfmethod}[Bootstrap resampling: resampling the data to estimate a sampling distribution, the sampling unit, and why i.i.d.\ resampling fails for dependent returns]\label{meth:bootstrap-resampling}
\qffield{Statement} The bootstrap draws subsamples ``with replacement'' from the data, computes the statistic on each, and repeats, typically 500 to 1{,}000 times; the results ``form an estimate of the sampling distribution of the statistic'', and percentiles of the sorted results give confidence intervals. ``The bootstrap is not appropriate for all distributions and statistics'': statistics driven by the extremes, such as the range, are poorly served \citep[sec.~1.3.3.4]{statistical-methods-handbook}. Applied to dependent data, the i.i.d.\ bootstrap ``scrambles'' the observations, so ``all dependence information is lost'' and it ``fails to give valid standard error estimates in the dependent case''; the block bootstrap instead joins blocks of consecutive observations chosen at random with replacement, and the block size must be chosen, a problem like bandwidth choice in smoothing \citep[secs.~2 and~2.3, pp.~221--223]{bootstrap-time-series-2003}.

\qffield{Notation} Data $x_1,\dots,x_n$; $B$ resamples; the \emph{sampling unit} is what is drawn: single observations (i.i.d.\ bootstrap), blocks of $b$ consecutive observations (moving-block bootstrap) or whole cross-sections of assets on a date (to keep cross-sectional correlation). Interpretation: as $B\to\infty$ the i.i.d.\ bootstrap variance of the mean tends to $\frac1n\cdot\frac1n\sum(x_i-\bar x)^2$, the \emph{ideal} bootstrap value; finite $B$ adds resampling noise to it.

\qffield{Conditions} The bootstrap approximates the sampling distribution by treating the sample as the population; it inherits the sample's size and biases and cannot create information on tails it did not observe. The i.i.d.\ version needs independent observations; for serially correlated returns, blocks (or other dependent-data schemes) are needed, and they remain biased in finite samples. Seeds and $B$ should be recorded so that resampling noise can be reproduced and judged.

\qffield{Practice} Examples, original. For the twelve active returns of \hyperref[meth:confidence-intervals]{confidence intervals}, 10{,}000 i.i.d.\ resamples give a standard error of the mean of 0.309 against the ideal 0.307, and a 2.5--97.5 percent interval of about $-0.09$ to 1.13, close to the $t$ interval. \emph{Dependence}: for 120 observations of an AR(1) series with coefficient 0.5, the exact variance of the sample mean is 0.0330. The ideal i.i.d.\ bootstrap variance has expectation 0.0108, a third of it; moving blocks of 12 recover 0.0261, 79 percent: better, still too low.

\qffield{Pitfalls} Resampling single days of autocorrelated or volatility-clustered returns and reporting the resulting narrow intervals. Resampling each asset separately and destroying cross-correlation. Choosing the block size, or reseeding, until a result appears.

\qffield{See also} \hyperref[meth:sequential-bootstrap]{sequential bootstrap} (overlapping labels); \hyperref[meth:resampled-efficiency]{resampled efficiency}; \hyperref[meth:sharpe-ratio-standard-error]{Sharpe-ratio standard errors}.

\end{qfmethod}

\begin{qfmethod}[Monte Carlo simulation: sampling from a specified model, the estimation target, seeds, standard error, the square-root rule, control variates and discretisation bias]\label{meth:monte-carlo-simulation}
\qffield{Statement} A Monte Carlo valuation samples paths of the variables from a specified model, computes the discounted payoff on each, and averages. With $\bar f$ the mean and $\omega$ the standard deviation of the discounted payoffs over $M$ trials, the standard error is $\omega/\sqrt M$ and a 95 percent confidence interval for the value is $\bar f\pm1.96\,\omega/\sqrt M$; uncertainty falls with the square root of the number of trials, so doubling the accuracy needs four times the trials \citep[chap.~21, sec.~21.6, pp.~492--494]{s7-hulloptions}. In Hull's example, 1{,}000 trials give an option value of 4.98 with a standard error of 0.24, and the closed-form price, 4.817, lies inside the interval (Example~21.8, p.~494).

\emph{Control variates.} When a similar derivative B has an analytic value $f_B$, both are simulated on the same random numbers and the estimate of A is corrected to $f_A=f_A^*-f_B^*+f_B$ (sec.~21.7, eq.~(21.20), p.~496). \emph{Discretisation against sampling error.} Simulating $\ln S$ with constant parameters is exact at any step, while an Euler step of $dS$ is exact only as $\Delta t\to0$ (p.~491 and n.~13). In general the time step contributes an error of order $\delta t$ and the finite number of paths one of order $N^{-1/2}$. For a fixed budget proportional to $N/\delta t$, Wilmott balances them with $N$ of order $K^{2/3}$ and $\delta t$ of order $K^{-1/3}$ \citep[chap.~80, secs.~80.4--80.6, pp.~1266--1268]{wide-wilmott-pwoqf}.

\qffield{Notation} Target: $\theta=E[g(Y)]$ under the specified distribution of $Y$, for example a discounted payoff in a risk-neutral model or a probability $P(Y\in A)=E[\mathbf 1_A(Y)]$, whose standard error is $\sqrt{\hat p(1-\hat p)/M}$. Seed: the starting state of the pseudo-random generator, which makes a run reproducible. Interpretation: with a control $X$ of known mean, $Y-b(X-E X)$ has the same mean as $Y$ and the smallest variance at $b^*=\mathrm{Cov}(Y,X)/\mathrm{Var}(X)$, where its variance is $(1-\rho^2)\mathrm{Var}(Y)$. Hull's correction is the case $b=1$. Estimating $b^*$ from the same draws adds a bias of order $1/M$, negligible at these sizes.

\qffield{Conditions} The estimate is only as good as the model: Monte Carlo gives the expectation under the assumed distribution, not under the true one. The standard-error formula assumes independent trials with finite payoff variance; dependent draws, such as overlapping paths, need other error estimates. Variance-reduction methods such as antithetic variables change the precision per trial, not the target (Hull, sec.~21.7). Seeds should be fixed before the run and reported, not chosen after seeing results. The standard error measures sampling error only: discretisation bias does not shrink with more paths, so it is measured separately, for example against the exact solution on common Brownian paths.

\qffield{Practice} Example, original. A stock at 100 follows geometric Brownian motion with volatility 20 percent; the risk-free rate is 3 percent; a one-year call struck at 105 has a closed-form value of 7.1281. With 100{,}000 draws the estimate is 7.1161 and the standard error 0.0396, so the closed form lies 0.3 standard errors away. The probability that the price ends below 90 is 0.2821 by simulation against 0.2820 exactly (standard error 0.0014). A quarter of the draws gives about twice the standard error, and antithetic pairs reduce it further.

With the discounted terminal price as control (known mean 100), the estimated coefficient is 0.542 and the correlation 0.875. The variance falls by $1/(1-\rho^2)=4.25$, and the estimate becomes 7.1289 with standard error 0.0192, against 0.0396 without the control. \emph{Discretisation}: with Euler steps of $dS$ on the same 400{,}000 Brownian paths as the exact solution, the value is biased by $-0.125$, $-0.054$ and $-0.025$ at 2, 4 and 8 steps. These biases are many times their standard errors (0.002 or less) and halve with the step, whereas more paths would leave them unchanged.

\qffield{Pitfalls} Reporting a Monte Carlo figure without its standard error, number of trials and seed. Reading agreement within sampling error as validation of the model. Comparing two simulated strategies on different random draws when common random numbers would sharpen the comparison.

\qffield{See also} \hyperref[def:financial-distributions]{distributions} (lognormal prices); \hyperref[mod:black-scholes-merton]{Black--Scholes--Merton}; \hyperref[meth:var-es-estimation-methods]{VaR and ES estimation} (Monte Carlo for risk); \hyperref[res:ito-lemma-geometric-brownian-motion]{It\^o's lemma} (exact and Euler sampling); \hyperref[meth:least-squares-monte-carlo]{least-squares Monte Carlo} (early exercise); \hyperref[meth:asian-option-valuation]{Asian options} (a geometric-average control).

\end{qfmethod}
